% Laget automatisk av elevark/lag_elevark.py fra opplegg/sentripetal-mobil.html. Endringer her blir overskrevet.
\documentclass[11pt]{article}
\usepackage{elevark}
\begin{document}
\elevarkhode{Fysikk 2 · Sirkelbevegelse}{Sentripetalakselerasjon med mobilen}{https://fysikklab.no/opplegg/sentripetal-mobil.html}{Opplegget på nett}

Mobilen har både et akselerometer og et gyroskop. Legger dere den i en salatslynge, kan phyphox måle akselerasjonen og vinkelfarten samtidig. Stemmer det at $a = ω^{2}r$, og kan dere finne ut hvor i mobilen akselerometeret sitter?

\section*{Hensikt}

\begin{itemize}
\item Måle sentripetalakselerasjonen ved ulike vinkelfarter.
\item Undersøke om $a$ er proporsjonal med $ω^{2}$, og bestemme radien $r$ fra stigningstallet.
\item Sammenligne $r$ med avstanden fra aksen til sensoren i mobilen.
\end{itemize}

\section*{Utstyr}

\begin{itemize}
\item Mobil med appen phyphox (gratis) og eksperimentet for sentripetalakselerasjon («Centripetal acceleration» på engelsk) under Mekanikk. Det viser akselerasjonen og vinkelfarten i samme graf.
\item Salatslynge med sveiv eller snor. En kontorstol går også, se under.
\item Et kjøkkenhåndkle til å kile fast mobilen, og en linjal.
\end{itemize}

\section*{Framgangsmåte}

\begin{enumerate}
\item Kil mobilen fast i salatslyngen med håndkleet. Mobilen skal ligge flatt med skjermen opp og ytterkanten mot veggen i slyngen.
\item Start målingen i phyphox, sett på lokket og sveiv jevnt. Hold farten konstant i 5–10 sekunder, og øk den så litt. Gjør 5–8 slike trinn, fra rolig til ganske raskt.
\item Stopp, ta ut mobilen og se på grafen over $a$ mot $ω$. Les av akselerasjonen og vinkelfarten midt i hvert trinn, og skriv dem inn i tabellen under. Dere kan også eksportere dataene fra phyphox.
\item Mål avstanden fra midten av slyngen (rotasjonsaksen) til midten av mobilen. Hvor sensoren sitter, finner dere ut i oppgave 4.
\end{enumerate}

\textbf{Uten salatslynge:} Sitt på en kontorstol og hold mobilen flatt på strak arm, eller tape den fast på setet nær kanten. La en annen snurre stolen rolig. Vinkelfarten blir mindre, men radien blir større.

\section*{Sikkerhet}

Sveiv rolig i starten, og sjekk at mobilen ligger godt fast. På kontorstolen skal dere snurre sakte, holde mobilen godt og stoppe hvis noen blir svimmel. Mange mobiler måler bare opp til omtrent 8~g eller 16~g. Det er nok i en vanlig salatslynge.

\section*{Oppgaver}

\begin{oppgaveliste}
\oppgave \fase{gjett}{Gjett}Hva skjer med akselerasjonen hvis dere sveiver dobbelt så fort? Hvor stor blir akselerasjonen 7~cm fra aksen når slyngen går rundt tre ganger per sekund? Sammenlign med $g$.
\svar
\oppgave \fase{test}{Mål}Gjør målingene og fyll inn tabellen. Hvordan ser grafen ut hvis dere plotter $a$ mot $ω$ i stedet for mot $ω^{2}$?
\svar
\oppgave \fase{test}{Regn}Ligger punktene på en rett linje gjennom origo når dere plotter $a$ mot $ω^{2}$? Hva er stigningstallet, og hvilken enhet har det?
\svar
\oppgave \fase{forklar}{Forklar}Stigningstallet er avstanden fra aksen til akselerometeret. Sammenlign med avstanden dere målte til midten av mobilen. Hvor i mobilen sitter sensoren?
\svar
\oppgave \fase{forklar}{Forklar}Hvilken kraft gir mobilen sentripetalakselerasjonen? Hvorfor føles det som om mobilen blir presset utover, selv om akselerasjonen peker innover?
\svar
\oppgave \fase{test}{Regn}Se på målingen med størst vinkelfart. Regn ut omløpstiden, banefarten og hvor stor kraft veggen må gi mobilen. Mobilen har massen 0,20~kg.
\svar
\oppgave \fase{utfordring}{Utfordring}Plott $a$ mot $v^{2}$, der $v = ωr$ er banefarten til sensoren. Hva blir stigningstallet nå? Hvorfor er det likevel lurere å plotte mot $ω^{2}$ i dette forsøket?
\svar
\end{oppgaveliste}

\end{document}
